\section{Mô tả chi tiết các ca sử dụng tiêu biểu}

%%%% UC-01: Đăng nhập %%%%
\subsection{Đăng nhập}
\subsubsection{Mô tả}
\begin{longtable}{|p{2.2cm}|p{13cm}|}
\caption{Đặc tả ca sử dụng Đăng nhập}
\label{tab:uc01} \\
\hline
\textbf{ID} & UC-01 \\
\hline
\textbf{Tên Use Case} & Đăng nhập \\
\hline
\textbf{Mô tả} & Cho phép người dùng đăng nhập vào hệ thống bằng tài khoản Google. \\
\hline
\textbf{Actor} & Đơn vị đào tạo, Học viên \\
\hline
\textbf{Tiền điều kiện} & 
\begin{minipage}[t]{\linewidth}
\begin{itemize}[nosep,after=\strut]
    \item Người dùng đã có tài khoản Google.
\end{itemize}
\end{minipage} \\
\hline
\textbf{Luồng chính} & 
\begin{minipage}[t]{\linewidth}
\vspace{2pt}
\begin{tabular}{p{6cm}|p{6cm}}
\textbf{Actor} & \textbf{System} \\
\hline
1. Mở trang đăng nhập. &  \\
 & 2. Hiển thị giao diện đăng nhập. \\
3. Chọn "Đăng nhập bằng Google" và tiến hành đăng nhập. & \\
 & 4. Kiểm tra thông tin đăng nhập từ Google. \\
 & 5. Nếu hợp lệ, tạo phiên đăng nhập (JWT token) và điều hướng người dùng vào giao diện tương ứng theo vai trò (Đơn vị đào tạo / Học viên). \\
\end{tabular}
\vspace{2pt}
\end{minipage} \\
\hline
\textbf{Luồng ngoại lệ} & 
\begin{minipage}[t]{\linewidth}
\begin{itemize}[nosep,after=\strut]
    \item \textbf{A1 -- Đăng nhập Google bị lỗi:}
    \begin{enumerate}[nosep]
        \item Hệ thống hiển thị thông báo lỗi.
        \item Người dùng quay lại bước 1 để đăng nhập lại.
    \end{enumerate}
\end{itemize}
\end{minipage} \\
\hline
\textbf{Hậu điều kiện} & 
\begin{minipage}[t]{\linewidth}
\begin{itemize}[nosep,after=\strut]
    \item Người dùng được xác thực và đăng nhập thành công.
    \item Phiên làm việc (JWT token) được tạo.
\end{itemize}
\end{minipage} \\
\hline
\textbf{Quy tắc nghiệp vụ} & 
\begin{minipage}[t]{\linewidth}
\begin{itemize}[nosep,after=\strut]
    \item Mỗi vai trò được điều hướng vào trang chức năng riêng.
    \item Phiên đăng nhập hết hạn sau 24 giờ.
\end{itemize}
\end{minipage} \\
\hline
\end{longtable}
Bảng \ref{tab:uc01} đặc tả ca sử dụng Đăng nhập.

\subsubsection{Biểu đồ tuần tự}
\begin{figure}[H]
    \centering
    \includegraphics[width=0.8\textwidth]{img/login-se.png}
    \caption{Biểu đồ tuần tự ca sử dụng Đăng nhập}
    \label{fig:uc01-se}
\end{figure}
Hình \ref{fig:uc01-se} trình bày về biểu đồ tuần tự ca sử dụng Đăng nhập.

\subsubsection{Biểu đồ hoạt động}
\begin{figure}[H]
    \centering
    \includegraphics[width=0.8\textwidth]{img/login-ac.png}
    \caption{Biểu đồ hoạt động ca sử dụng Đăng nhập}
    \label{fig:uc01-ac}
\end{figure}
Hình \ref{fig:uc01-ac} trình bày về biểu đồ hoạt động ca sử dụng Đăng nhập.

%%%% UC-02: Thêm khóa học %%%%
\subsection{Thêm khóa học}
\subsubsection{Mô tả}
\begin{longtable}{|p{2.2cm}|p{13cm}|}
\caption{Đặc tả ca sử dụng Thêm khóa học}
\label{tab:uc02} \\
\hline
\textbf{ID} & UC-02 \\
\hline
\textbf{Tên Use Case} & Thêm khóa học \\
\hline
\textbf{Mô tả} & Cho phép Đơn vị đào tạo thêm mới một khóa học vào hệ thống. \\
\hline
\textbf{Actor} & Đơn vị đào tạo \\
\hline
\textbf{Tiền điều kiện} & 
\begin{minipage}[t]{\linewidth}
\begin{itemize}[nosep,after=\strut]
    \item Đơn vị đào tạo đã đăng nhập vào hệ thống.
\end{itemize}
\end{minipage} \\
\hline
\textbf{Luồng chính} & 
\begin{minipage}[t]{\linewidth}
\vspace{2pt}
\begin{tabular}{p{6cm}|p{6cm}}
\textbf{Actor} & \textbf{System} \\
\hline
1. Truy cập trang "Khóa học". & \\
 & 2. Hiển thị danh sách khóa học hiện có. \\
3. Chọn "Thêm khóa học". & \\
 & 4. Hiển thị biểu mẫu thêm khóa học. \\
5. Nhập thông tin: tên khóa học, mô tả, thời lượng, nội dung đào tạo. & \\
6. Nhấn "Lưu". & \\
 & 7. Kiểm tra tính hợp lệ của dữ liệu. \\
 & 8. Lưu khóa học vào cơ sở dữ liệu và thông báo thành công. \\
\end{tabular}
\vspace{2pt}
\end{minipage} \\
\hline
\textbf{Luồng ngoại lệ} & 
\begin{minipage}[t]{\linewidth}
\begin{itemize}[nosep,after=\strut]
    \item \textbf{A1 -- Thiếu thông tin bắt buộc:} Hệ thống hiển thị lỗi và yêu cầu điền đầy đủ thông tin.
    \item \textbf{A2 -- Tên khóa học bị trùng:} Hệ thống từ chối thêm mới và thông báo lỗi.
\end{itemize}
\end{minipage} \\
\hline
\textbf{Hậu điều kiện} & 
\begin{minipage}[t]{\linewidth}
\begin{itemize}[nosep,after=\strut]
    \item Khóa học được thêm mới vào cơ sở dữ liệu.
\end{itemize}
\end{minipage} \\
\hline
\textbf{Quy tắc nghiệp vụ} & 
\begin{minipage}[t]{\linewidth}
\begin{itemize}[nosep,after=\strut]
    \item Mã khóa học phải là duy nhất.
\end{itemize}
\end{minipage} \\
\hline
\end{longtable}
Bảng \ref{tab:uc02} đặc tả ca sử dụng Thêm khóa học.

\subsubsection{Biểu đồ tuần tự}
\begin{figure}[H]
    \centering
    \includegraphics[width=0.8\textwidth]{img/add-course-se.png}
    \caption{Biểu đồ tuần tự ca sử dụng Thêm khóa học}
    \label{fig:uc02-se}
\end{figure}
Hình \ref{fig:uc02-se} trình bày về biểu đồ tuần tự ca sử dụng Thêm khóa học.

\subsubsection{Biểu đồ hoạt động}
\begin{figure}[H]
    \centering
    \includegraphics[width=0.8\textwidth]{img/add-course-ac.png}
    \caption{Biểu đồ hoạt động ca sử dụng Thêm khóa học}
    \label{fig:uc02-ac}
\end{figure}
Hình \ref{fig:uc02-ac} trình bày về biểu đồ hoạt động ca sử dụng Thêm khóa học.

%%%% UC-03: Sửa khóa học %%%%
\subsection{Sửa khóa học}
\subsubsection{Mô tả}
\begin{longtable}{|p{2.2cm}|p{13cm}|}
\caption{Đặc tả ca sử dụng Sửa khóa học}
\label{tab:uc03} \\
\hline
\textbf{ID} & UC-03 \\
\hline
\textbf{Tên Use Case} & Sửa khóa học \\
\hline
\textbf{Mô tả} & Cho phép Đơn vị đào tạo chỉnh sửa thông tin của một khóa học đã có trong hệ thống. \\
\hline
\textbf{Actor} & Đơn vị đào tạo \\
\hline
\textbf{Tiền điều kiện} & 
\begin{minipage}[t]{\linewidth}
\begin{itemize}[nosep,after=\strut]
    \item Đơn vị đào tạo đã đăng nhập vào hệ thống.
    \item Khóa học cần sửa đã tồn tại trong hệ thống.
\end{itemize}
\end{minipage} \\
\hline
\textbf{Luồng chính} & 
\begin{minipage}[t]{\linewidth}
\vspace{2pt}
\begin{tabular}{p{6cm}|p{6cm}}
\textbf{Actor} & \textbf{System} \\
\hline
1. Truy cập trang "Khóa học". & \\
 & 2. Hiển thị danh sách khóa học hiện có. \\
3. Chọn khóa học cần sửa. & \\
 & 4. Hiển thị thông tin chi tiết khóa học. \\
5. Chỉnh sửa các thông tin cần thay đổi. & \\
6. Nhấn "Lưu". & \\
 & 7. Kiểm tra tính hợp lệ của dữ liệu. \\
 & 8. Cập nhật thông tin khóa học vào cơ sở dữ liệu và thông báo thành công. \\
\end{tabular}
\vspace{2pt}
\end{minipage} \\
\hline
\textbf{Luồng ngoại lệ} & 
\begin{minipage}[t]{\linewidth}
\begin{itemize}[nosep,after=\strut]
    \item \textbf{A1 -- Thiếu thông tin bắt buộc:} Hệ thống hiển thị lỗi và yêu cầu điền đầy đủ thông tin.
    \item \textbf{A2 -- Tên khóa học bị trùng:} Hệ thống từ chối cập nhật và thông báo lỗi.
\end{itemize}
\end{minipage} \\
\hline
\textbf{Hậu điều kiện} & 
\begin{minipage}[t]{\linewidth}
\begin{itemize}[nosep,after=\strut]
    \item Thông tin khóa học được cập nhật thành công trong cơ sở dữ liệu.
\end{itemize}
\end{minipage} \\
\hline
\textbf{Quy tắc nghiệp vụ} & 
\begin{minipage}[t]{\linewidth}
\begin{itemize}[nosep,after=\strut]
    \item Mã khóa học phải là duy nhất.
\end{itemize}
\end{minipage} \\
\hline
\end{longtable}
Bảng \ref{tab:uc03} đặc tả ca sử dụng Sửa khóa học.

\subsubsection{Biểu đồ tuần tự}
\begin{figure}[H]
    \centering
    \includegraphics[width=0.8\textwidth]{img/edit-course-se.png}
    \caption{Biểu đồ tuần tự ca sử dụng Sửa khóa học}
    \label{fig:uc03-se}
\end{figure}
Hình \ref{fig:uc03-se} trình bày về biểu đồ tuần tự ca sử dụng Sửa khóa học.

\subsubsection{Biểu đồ hoạt động}
\begin{figure}[H]
    \centering
    \includegraphics[width=0.8\textwidth]{img/edit-course-ac.png}
    \caption{Biểu đồ hoạt động ca sử dụng Sửa khóa học}
    \label{fig:uc03-ac}
\end{figure}
Hình \ref{fig:uc03-ac} trình bày về biểu đồ hoạt động ca sử dụng Sửa khóa học.


%%%% UC-04: Xóa khóa học %%%%
\subsection{Xóa khóa học}
\subsubsection{Mô tả}
\begin{longtable}{|p{2.2cm}|p{13cm}|}
\caption{Đặc tả ca sử dụng Xóa khóa học}
\label{tab:uc04} \\
\hline
\textbf{ID} & UC-04 \\
\hline
\textbf{Tên Use Case} & Xóa khóa học \\
\hline
\textbf{Mô tả} & Cho phép Đơn vị đào tạo xóa một khóa học khỏi hệ thống. \\
\hline
\textbf{Actor} & Đơn vị đào tạo \\
\hline
\textbf{Tiền điều kiện} & 
\begin{minipage}[t]{\linewidth}
\begin{itemize}[nosep,after=\strut]
    \item Đơn vị đào tạo đã đăng nhập vào hệ thống.
    \item Khóa học cần xóa đã tồn tại trong hệ thống.
\end{itemize}
\end{minipage} \\
\hline
\textbf{Luồng chính} & 
\begin{minipage}[t]{\linewidth}
\vspace{2pt}
\begin{tabular}{p{6cm}|p{6cm}}
\textbf{Actor} & \textbf{System} \\
\hline
1. Truy cập trang "Khóa học". & \\
 & 2. Hiển thị danh sách khóa học hiện có. \\
3. Chọn khóa học cần xóa. & \\
4. Nhấn "Xóa" và xác nhận thao tác. & \\
 & 5. Kiểm tra ràng buộc dữ liệu (khóa học có đang gắn với chứng chỉ không). \\
 & 6. Xóa khóa học khỏi cơ sở dữ liệu và thông báo thành công. \\
\end{tabular}
\vspace{2pt}
\end{minipage} \\
\hline
\textbf{Luồng ngoại lệ} & 
\begin{minipage}[t]{\linewidth}
\begin{itemize}[nosep,after=\strut]
    \item \textbf{A1 -- Khóa học đang gắn với chứng chỉ:} Hệ thống từ chối thao tác xóa và thông báo lý do.
\end{itemize}
\end{minipage} \\
\hline
\textbf{Hậu điều kiện} & 
\begin{minipage}[t]{\linewidth}
\begin{itemize}[nosep,after=\strut]
    \item Khóa học được xóa khỏi cơ sở dữ liệu.
\end{itemize}
\end{minipage} \\
\hline
\textbf{Quy tắc nghiệp vụ} & 
\begin{minipage}[t]{\linewidth}
\begin{itemize}[nosep,after=\strut]
    \item Không được xóa khóa học đang được gắn với chứng chỉ.
\end{itemize}
\end{minipage} \\
\hline
\end{longtable}
Bảng \ref{tab:uc04} đặc tả ca sử dụng Xóa khóa học.

\subsubsection{Biểu đồ tuần tự}
\begin{figure}[H]
    \centering
    \includegraphics[width=0.8\textwidth]{img/delete-course-se.png}
    \caption{Biểu đồ tuần tự ca sử dụng Xóa khóa học}
    \label{fig:uc04-se}
\end{figure}
Hình \ref{fig:uc04-se} trình bày về biểu đồ tuần tự ca sử dụng Xóa khóa học.

\subsubsection{Biểu đồ hoạt động}
\begin{figure}[H]
    \centering
    \includegraphics[width=0.8\textwidth]{img/delete-course-ac.png}
    \caption{Biểu đồ hoạt động ca sử dụng Xóa khóa học}
    \label{fig:uc04-ac}
\end{figure}
Hình \ref{fig:uc04-ac} trình bày về biểu đồ hoạt động ca sử dụng Xóa khóa học.


%%%% UC-05: Thêm học viên %%%%
\subsection{Thêm học viên}
\subsubsection{Mô tả}
\begin{longtable}{|p{2.2cm}|p{13cm}|}
\caption{Đặc tả ca sử dụng Thêm học viên}
\label{tab:uc05} \\
\hline
\textbf{ID} & UC-05 \\
\hline
\textbf{Tên Use Case} & Thêm học viên \\
\hline
\textbf{Mô tả} & Cho phép Đơn vị đào tạo thêm mới một học viên vào hệ thống. \\
\hline
\textbf{Actor} & Đơn vị đào tạo \\
\hline
\textbf{Tiền điều kiện} & 
\begin{minipage}[t]{\linewidth}
\begin{itemize}[nosep,after=\strut]
    \item Đơn vị đào tạo đã đăng nhập vào hệ thống.
\end{itemize}
\end{minipage} \\
\hline
\textbf{Luồng chính} & 
\begin{minipage}[t]{\linewidth}
\vspace{2pt}
\begin{tabular}{p{6cm}|p{6cm}}
\textbf{Actor} & \textbf{System} \\
\hline
1. Truy cập trang "Quản lý học viên". & \\
 & 2. Hiển thị danh sách học viên hiện có. \\
3. Chọn "Thêm học viên". & \\
 & 4. Hiển thị biểu mẫu thêm học viên. \\
5. Nhập thông tin: mã học viên, tên học viên, email. & \\
6. Nhấn "Lưu". & \\
 & 7. Kiểm tra tính hợp lệ của dữ liệu. \\
 & 8. Lưu học viên vào cơ sở dữ liệu và thông báo thành công. \\
\end{tabular}
\vspace{2pt}
\end{minipage} \\
\hline
\textbf{Luồng ngoại lệ} & 
\begin{minipage}[t]{\linewidth}
\begin{itemize}[nosep,after=\strut]
    \item \textbf{A1 -- Thiếu thông tin bắt buộc:} Hệ thống hiển thị lỗi và yêu cầu điền đầy đủ thông tin.
    \item \textbf{A2 -- Mã học viên bị trùng:} Hệ thống từ chối thêm mới và thông báo lỗi.
\end{itemize}
\end{minipage} \\
\hline
\textbf{Hậu điều kiện} & 
\begin{minipage}[t]{\linewidth}
\begin{itemize}[nosep,after=\strut]
    \item Học viên được thêm mới vào cơ sở dữ liệu.
\end{itemize}
\end{minipage} \\
\hline
\textbf{Quy tắc nghiệp vụ} & 
\begin{minipage}[t]{\linewidth}
\begin{itemize}[nosep,after=\strut]
    \item Mã học viên phải là duy nhất.
\end{itemize}
\end{minipage} \\
\hline
\end{longtable}
Bảng \ref{tab:uc05} đặc tả ca sử dụng Thêm học viên.

\subsubsection{Biểu đồ tuần tự}
\begin{figure}[H]
    \centering
    \includegraphics[width=0.8\textwidth]{img/add-student-se.png}
    \caption{Biểu đồ tuần tự ca sử dụng Thêm học viên}
    \label{fig:uc05-se}
\end{figure}
Hình \ref{fig:uc05-se} trình bày về biểu đồ tuần tự ca sử dụng Thêm học viên.

\subsubsection{Biểu đồ hoạt động}
\begin{figure}[H]
    \centering
    \includegraphics[width=0.8\textwidth]{img/add-student-ac.png}
    \caption{Biểu đồ hoạt động ca sử dụng Thêm học viên}
    \label{fig:uc05-ac}
\end{figure}
Hình \ref{fig:uc05-ac} trình bày về biểu đồ hoạt động ca sử dụng Thêm học viên.


%%%% UC-06: Sửa học viên %%%%
\subsection{Sửa học viên}
\subsubsection{Mô tả}
\begin{longtable}{|p{2.2cm}|p{13cm}|}
\caption{Đặc tả ca sử dụng Sửa học viên}
\label{tab:uc06} \\
\hline
\textbf{ID} & UC-06 \\
\hline
\textbf{Tên Use Case} & Sửa học viên \\
\hline
\textbf{Mô tả} & Cho phép Đơn vị đào tạo chỉnh sửa thông tin của một học viên đã có trong hệ thống. \\
\hline
\textbf{Actor} & Đơn vị đào tạo \\
\hline
\textbf{Tiền điều kiện} & 
\begin{minipage}[t]{\linewidth}
\begin{itemize}[nosep,after=\strut]
    \item Đơn vị đào tạo đã đăng nhập vào hệ thống.
    \item Học viên cần sửa đã tồn tại trong hệ thống.
\end{itemize}
\end{minipage} \\
\hline
\textbf{Luồng chính} & 
\begin{minipage}[t]{\linewidth}
\vspace{2pt}
\begin{tabular}{p{6cm}|p{6cm}}
\textbf{Actor} & \textbf{System} \\
\hline
1. Truy cập trang "Quản lý học viên". & \\
 & 2. Hiển thị danh sách học viên hiện có. \\
3. Chọn học viên cần sửa. & \\
 & 4. Hiển thị thông tin chi tiết học viên. \\
5. Chỉnh sửa các thông tin cần thay đổi. & \\
6. Nhấn "Lưu". & \\
 & 7. Kiểm tra tính hợp lệ của dữ liệu. \\
 & 8. Cập nhật thông tin học viên vào cơ sở dữ liệu và thông báo thành công. \\
\end{tabular}
\vspace{2pt}
\end{minipage} \\
\hline
\textbf{Luồng ngoại lệ} & 
\begin{minipage}[t]{\linewidth}
\begin{itemize}[nosep,after=\strut]
    \item \textbf{A1 -- Thiếu thông tin bắt buộc:} Hệ thống hiển thị lỗi và yêu cầu điền đầy đủ thông tin.
\end{itemize}
\end{minipage} \\
\hline
\textbf{Hậu điều kiện} & 
\begin{minipage}[t]{\linewidth}
\begin{itemize}[nosep,after=\strut]
    \item Thông tin học viên được cập nhật thành công trong cơ sở dữ liệu.
\end{itemize}
\end{minipage} \\
\hline
\textbf{Quy tắc nghiệp vụ} & 
\begin{minipage}[t]{\linewidth}
\begin{itemize}[nosep,after=\strut]
    \item Mã học viên phải là duy nhất.
\end{itemize}
\end{minipage} \\
\hline
\end{longtable}
Bảng \ref{tab:uc06} đặc tả ca sử dụng Sửa học viên.

\subsubsection{Biểu đồ tuần tự}
\begin{figure}[H]
    \centering
    \includegraphics[width=0.8\textwidth]{img/edit-student-se.png}
    \caption{Biểu đồ tuần tự ca sử dụng Sửa học viên}
    \label{fig:uc06-se}
\end{figure}
Hình \ref{fig:uc06-se} trình bày về biểu đồ tuần tự ca sử dụng Sửa học viên.

\subsubsection{Biểu đồ hoạt động}
\begin{figure}[H]
    \centering
    \includegraphics[width=0.8\textwidth]{img/edit-student-ac.png}
    \caption{Biểu đồ hoạt động ca sử dụng Sửa học viên}
    \label{fig:uc06-ac}
\end{figure}
Hình \ref{fig:uc06-ac} trình bày về biểu đồ hoạt động ca sử dụng Sửa học viên.


%%%% UC-07: Xóa học viên %%%%
\subsection{Xóa học viên}
\subsubsection{Mô tả}
\begin{longtable}{|p{2.2cm}|p{13cm}|}
\caption{Đặc tả ca sử dụng Xóa học viên}
\label{tab:uc07} \\
\hline
\textbf{ID} & UC-07 \\
\hline
\textbf{Tên Use Case} & Xóa học viên \\
\hline
\textbf{Mô tả} & Cho phép Đơn vị đào tạo xóa một học viên khỏi hệ thống. \\
\hline
\textbf{Actor} & Đơn vị đào tạo \\
\hline
\textbf{Tiền điều kiện} & 
\begin{minipage}[t]{\linewidth}
\begin{itemize}[nosep,after=\strut]
    \item Đơn vị đào tạo đã đăng nhập vào hệ thống.
    \item Học viên cần xóa đã tồn tại trong hệ thống.
\end{itemize}
\end{minipage} \\
\hline
\textbf{Luồng chính} & 
\begin{minipage}[t]{\linewidth}
\vspace{2pt}
\begin{tabular}{p{6cm}|p{6cm}}
\textbf{Actor} & \textbf{System} \\
\hline
1. Truy cập trang "Quản lý học viên". & \\
 & 2. Hiển thị danh sách học viên hiện có. \\
3. Chọn học viên cần xóa. & \\
4. Nhấn "Xóa" và xác nhận thao tác. & \\
 & 5. Kiểm tra ràng buộc dữ liệu (học viên có chứng chỉ không). \\
 & 6. Xóa học viên khỏi cơ sở dữ liệu và thông báo thành công. \\
\end{tabular}
\vspace{2pt}
\end{minipage} \\
\hline
\textbf{Luồng ngoại lệ} & 
\begin{minipage}[t]{\linewidth}
\begin{itemize}[nosep,after=\strut]
    \item \textbf{A1 -- Học viên đã có chứng chỉ:} Hệ thống từ chối thao tác xóa và thông báo lý do.
\end{itemize}
\end{minipage} \\
\hline
\textbf{Hậu điều kiện} & 
\begin{minipage}[t]{\linewidth}
\begin{itemize}[nosep,after=\strut]
    \item Học viên được xóa khỏi cơ sở dữ liệu.
\end{itemize}
\end{minipage} \\
\hline
\textbf{Quy tắc nghiệp vụ} & 
\begin{minipage}[t]{\linewidth}
\begin{itemize}[nosep,after=\strut]
    \item Không được xóa học viên đã có chứng chỉ.
\end{itemize}
\end{minipage} \\
\hline
\end{longtable}
Bảng \ref{tab:uc07} đặc tả ca sử dụng Xóa học viên.

\subsubsection{Biểu đồ tuần tự}
\begin{figure}[H]
    \centering
    \includegraphics[width=0.8\textwidth]{img/delete-student-se.png}
    \caption{Biểu đồ tuần tự ca sử dụng Xóa học viên}
    \label{fig:uc07-se}
\end{figure}
Hình \ref{fig:uc07-se} trình bày về biểu đồ tuần tự ca sử dụng Xóa học viên.

\subsubsection{Biểu đồ hoạt động}
\begin{figure}[H]
    \centering
    \includegraphics[width=0.8\textwidth]{img/delete-student-ac.png}
    \caption{Biểu đồ hoạt động ca sử dụng Xóa học viên}
    \label{fig:uc07-ac}
\end{figure}
Hình \ref{fig:uc07-ac} trình bày về biểu đồ hoạt động ca sử dụng Xóa học viên.


%%%% UC-08: Cấp chứng chỉ %%%%
\subsection{Cấp chứng chỉ}
\subsubsection{Mô tả}
\begin{longtable}{|p{2.2cm}|p{13cm}|}
\caption{Đặc tả ca sử dụng Cấp chứng chỉ}
\label{tab:uc08} \\
\hline
\textbf{ID} & UC-08 \\
\hline
\textbf{Tên Use Case} & Cấp chứng chỉ \\
\hline
\textbf{Mô tả} & Cho phép Đơn vị đào tạo cấp mới một chứng chỉ NFT cho học viên. \\
\hline
\textbf{Actor} & Đơn vị đào tạo \\
\hline
\textbf{Tiền điều kiện} & 
\begin{minipage}[t]{\linewidth}
\begin{itemize}[nosep,after=\strut]
    \item Đơn vị đào tạo đã đăng nhập vào hệ thống.
    \item Học viên đã tồn tại trong hệ thống.
\end{itemize}
\end{minipage} \\
\hline
\textbf{Luồng chính} & 
\begin{minipage}[t]{\linewidth}
\vspace{2pt}
\begin{tabular}{p{6cm}|p{6cm}}
\textbf{Actor} & \textbf{System} \\
\hline
1. Truy cập trang "Chứng chỉ". & \\
 & 2. Hiển thị danh sách chứng chỉ đã cấp. \\
3. Chọn "Cấp chứng chỉ mới". & \\
 & 4. Hiển thị biểu mẫu cấp chứng chỉ. \\
5. Nhập thông tin: mã học viên, khóa học, ngày cấp, ngày hết hạn. & \\
6. Nhấn "Cấp chứng chỉ". & \\
 & 7. Kiểm tra tính hợp lệ của dữ liệu. \\
 & 8. Ghi giao dịch lên blockchain và lưu chứng chỉ vào cơ sở dữ liệu với trạng thái "Issued". \\
 & 9. Thông báo cấp chứng chỉ thành công. \\
\end{tabular}
\vspace{2pt}
\end{minipage} \\

\hline
\textbf{Luồng ngoại lệ} & 
\begin{minipage}[t]{\linewidth}
\begin{itemize}[nosep,after=\strut]
    \item \textbf{A1 -- Thiếu thông tin bắt buộc:} Hệ thống hiển thị lỗi và yêu cầu điền đầy đủ thông tin.
    \item \textbf{A2 -- Học viên không tồn tại:} Hệ thống từ chối cấp chứng chỉ và thông báo lỗi.
    \item \textbf{A3 -- Lỗi kết nối blockchain hoặc lỗi ghi giao dịch:} Hệ thống thông báo thao tác không thành công.
\end{itemize}
\end{minipage} \\
\hline
\textbf{Hậu điều kiện} & 
\begin{minipage}[t]{\linewidth}
\begin{itemize}[nosep,after=\strut]
    \item Chứng chỉ được cấp mới với trạng thái "Issued" và lưu vào hệ thống.
\end{itemize}
\end{minipage} \\
\hline
\textbf{Quy tắc nghiệp vụ} & 
\begin{minipage}[t]{\linewidth}
\begin{itemize}[nosep,after=\strut]
    \item Mã xác minh chứng chỉ phải là duy nhất.
    \item Ngày hết hạn phải lớn hơn ngày hiện tại.
\end{itemize}
\end{minipage} \\
\hline
\end{longtable}
Bảng \ref{tab:uc08} đặc tả ca sử dụng Cấp chứng chỉ.

\subsubsection{Biểu đồ tuần tự}
\begin{figure}[H]
    \centering
    \includegraphics[width=0.8\textwidth]{img/issue-cert-se.png}
    \caption{Biểu đồ tuần tự ca sử dụng Cấp chứng chỉ}
    \label{fig:uc08-se}
\end{figure}
Hình \ref{fig:uc08-se} trình bày về biểu đồ tuần tự ca sử dụng Cấp chứng chỉ.

\subsubsection{Biểu đồ hoạt động}
\begin{figure}[H]
    \centering
    \includegraphics[width=0.38\textwidth]{img/issue-cert-ac.png}
    \caption{Biểu đồ hoạt động ca sử dụng Cấp chứng chỉ}
    \label{fig:uc08-ac}
\end{figure}
Hình \ref{fig:uc08-ac} trình bày về biểu đồ hoạt động ca sử dụng Cấp chứng chỉ.


%%%% UC-09: Thay thế chứng chỉ %%%%
\subsection{Thay thế chứng chỉ}
\subsubsection{Mô tả}
\begin{longtable}{|p{2.2cm}|p{13cm}|}
\caption{Đặc tả ca sử dụng Thay thế chứng chỉ}
\label{tab:uc09} \\
\hline
\textbf{ID} & UC-09 \\
\hline
\textbf{Tên Use Case} & Thay thế chứng chỉ \\
\hline
\textbf{Mô tả} & Cho phép Đơn vị đào tạo thay thế một chứng chỉ NFT đã cấp (do sai thông tin, cập nhật dữ liệu…) bằng chứng chỉ mới. \\
\hline
\textbf{Actor} & Đơn vị đào tạo \\
\hline
\textbf{Tiền điều kiện} & 
\begin{minipage}[t]{\linewidth}
\begin{itemize}[nosep,after=\strut]
    \item Đơn vị đào tạo đã đăng nhập vào hệ thống.
    \item Chứng chỉ cần thay thế đã tồn tại trong hệ thống.
\end{itemize}
\end{minipage} \\
\hline
\textbf{Luồng chính} & 
\begin{minipage}[t]{\linewidth}
\vspace{2pt}
\begin{tabular}{p{6cm}|p{6cm}}
\textbf{Actor} & \textbf{System} \\
\hline
1. Truy cập trang "Chứng chỉ". & \\
 & 2. Hiển thị danh sách chứng chỉ đã cấp. \\
3. Chọn chứng chỉ cần thay thế. & \\
 & 4. Hiển thị thông tin chi tiết chứng chỉ. \\
5. Chọn "Thay thế" và cập nhật thôngময় tin mới. & \\
6. Nhấn "Xác nhận thay thế". & \\
 & 7. Kiểm tra tính hợp lệ của dữ liệu. \\
 & 8. Đánh dấu chứng chỉ cũ là "Replaced", tạo chứng chỉ mới trên blockchain. \\
 & 9. Thông báo thay thế chứng chỉ thành công. \\
\end{tabular}
\vspace{2pt}
\end{minipage} \\
\hline
\textbf{Luồng ngoại lệ} & 
\begin{minipage}[t]{\linewidth}
\begin{itemize}[nosep,after=\strut]
    \item \textbf{A1 -- Thiếu thông tin bắt buộc:} Hệ thống hiển thị lỗi và yêu cầu điền đầy đủ thông tin.
    \item \textbf{A2 -- Lỗi kết nối blockchain hoặc lỗi ghi giao dịch:} Hệ thống thông báo thao tác không thành công.
\end{itemize}
\end{minipage} \\
\hline
\textbf{Hậu điều kiện} & 
\begin{minipage}[t]{\linewidth}
\begin{itemize}[nosep,after=\strut]
    \item Chứng chỉ cũ được đánh dấu "Replaced".
    \item Chứng chỉ mới được tạo với trạng thái "Issued".
\end{itemize}
\end{minipage} \\
\hline
\textbf{Quy tắc nghiệp vụ} & 
\begin{minipage}[t]{\linewidth}
\begin{itemize}[nosep,after=\strut]
    \item Chỉ người cấp (issuer) mới có quyền thay thế chứng chỉ do mình cấp.
\end{itemize}
\end{minipage} \\
\hline
\end{longtable}
Bảng \ref{tab:uc09} đặc tả ca sử dụng Thay thế chứng chỉ.

\subsubsection{Biểu đồ tuần tự}
\begin{figure}[H]
    \centering
    \includegraphics[width=0.8\textwidth]{img/replace-cert-se.png}
    \caption{Biểu đồ tuần tự ca sử dụng Thay thế chứng chỉ}
    \label{fig:uc09-se}
\end{figure}
Hình \ref{fig:uc09-se} trình bày về biểu đồ tuần tự ca sử dụng Thay thế chứng chỉ.

\subsubsection{Biểu đồ hoạt động}
\begin{figure}[H]
    \centering
    \includegraphics[width=0.6\textwidth]{img/replace-cert-ac.png}
    \caption{Biểu đồ hoạt động ca sử dụng Thay thế chứng chỉ}
    \label{fig:uc09-ac}
\end{figure}
Hình \ref{fig:uc09-ac} trình bày về biểu đồ hoạt động ca sử dụng Thay thế chứng chỉ.


%%%% UC-10: Thu hồi chứng chỉ %%%%
\subsection{Thu hồi chứng chỉ}
\subsubsection{Mô tả}
\begin{longtable}{|p{2.2cm}|p{13cm}|}
\caption{Đặc tả ca sử dụng Thu hồi chứng chỉ}
\label{tab:uc10} \\
\hline
\textbf{ID} & UC-10 \\
\hline
\textbf{Tên Use Case} & Thu hồi chứng chỉ \\
\hline
\textbf{Mô tả} & Cho phép Đơn vị đào tạo thu hồi một chứng chỉ NFT đã cấp cho học viên. \\
\hline
\textbf{Actor} & Đơn vị đào tạo \\
\hline
\textbf{Tiền điều kiện} & 
\begin{minipage}[t]{\linewidth}
\begin{itemize}[nosep,after=\strut]
    \item Đơn vị đào tạo đã đăng nhập vào hệ thống.
    \item Chứng chỉ cần thu hồi đã tồn tại và đang ở trạng thái "Issued" hoặc "Active".
\end{itemize}
\end{minipage} \\
\hline
\textbf{Luồng chính} & 
\begin{minipage}[t]{\linewidth}
\vspace{2pt}
\begin{tabular}{p{6cm}|p{6cm}}
\textbf{Actor} & \textbf{System} \\
\hline
1. Truy cập trang "Chứng chỉ". & \\
 & 2. Hiển thị danh sách chứng chỉ đã cấp. \\
3. Chọn chứng chỉ cần thu hồi. & \\
 & 4. Hiển thị thông tin chi tiết chứng chỉ. \\
5. Chọn "Thu hồi" và nhập lý do thu hồi. & \\
6. Nhấn "Xác nhận thu hồi". & \\
 & 7. Cập nhật trạng thái chứng chỉ thành "Revoked" trên blockchain. \\
 & 8. Thông báo thu hồi chứng chỉ thành công. \\
\end{tabular}
\vspace{2pt}
\end{minipage} \\
\hline
\textbf{Luồng ngoại lệ} & 
\begin{minipage}[t]{\linewidth}
\begin{itemize}[nosep,after=\strut]
    \item \textbf{A1 -- Lỗi kết nối blockchain hoặc lỗi ghi giao dịch:} Hệ thống thông báo thao tác không thành công.
\end{itemize}
\end{minipage} \\
\hline
\textbf{Hậu điều kiện} & 
\begin{minipage}[t]{\linewidth}
\begin{itemize}[nosep,after=\strut]
    \item Chứng chỉ được đánh dấu "Revoked" trên blockchain và cơ sở dữ liệu.
\end{itemize}
\end{minipage} \\
\hline
\textbf{Quy tắc nghiệp vụ} & 
\begin{minipage}[t]{\linewidth}
\begin{itemize}[nosep,after=\strut]
    \item Không thể phục hồi chứng chỉ đã thu hồi.
    \item Chỉ người cấp (issuer) mới có quyền thu hồi chứng chỉ do mình cấp.
\end{itemize}
\end{minipage} \\
\hline
\end{longtable}
Bảng \ref{tab:uc10} đặc tả ca sử dụng Thu hồi chứng chỉ.

\subsubsection{Biểu đồ tuần tự}
\begin{figure}[H]
    \centering
    \includegraphics[width=0.8\textwidth]{img/revoke-cert-se.png}
    \caption{Biểu đồ tuần tự ca sử dụng Thu hồi chứng chỉ}
    \label{fig:uc10-se}
\end{figure}
Hình \ref{fig:uc10-se} trình bày về biểu đồ tuần tự ca sử dụng Thu hồi chứng chỉ.

\subsubsection{Biểu đồ hoạt động}
\begin{figure}[H]
    \centering
    \includegraphics[width=0.8\textwidth]{img/revoke-cert-ac.png}
    \caption{Biểu đồ hoạt động ca sử dụng Thu hồi chứng chỉ}
    \label{fig:uc10-ac}
\end{figure}
Hình \ref{fig:uc10-ac} trình bày về biểu đồ hoạt động ca sử dụng Thu hồi chứng chỉ.


%%%% UC-11: Cập nhật địa chỉ ví %%%%
\subsection{Cập nhật địa chỉ ví}
\subsubsection{Mô tả}
\begin{longtable}{|p{2.2cm}|p{13cm}|}
\caption{Đặc tả ca sử dụng Cập nhật địa chỉ ví}
\label{tab:uc11} \\
\hline
\textbf{ID} & UC-11 \\
\hline
\textbf{Tên Use Case} & Cập nhật địa chỉ ví \\
\hline
\textbf{Mô tả} & Cho phép học viên cập nhật hoặc thay đổi địa chỉ ví blockchain (Ethereum wallet address) được sử dụng để nhận chứng chỉ NFT trên hệ thống. \\
\hline
\textbf{Actor} & Học viên \\
\hline
\textbf{Tiền điều kiện} & 
\begin{minipage}[t]{\linewidth}
\begin{itemize}[nosep,after=\strut]
    \item Học viên đã đăng nhập vào hệ thống.
    \item Học viên có ví Metamask.
    \item Địa chỉ ví phải là ví do học viên sở hữu.
\end{itemize}
\end{minipage} \\
\hline
\textbf{Luồng chính} & 
\begin{minipage}[t]{\linewidth}
\vspace{2pt}
\begin{tabular}{p{6cm}|p{6cm}}
\textbf{Actor} & \textbf{System} \\
\hline
1. Truy cập mục "Ví Blockchain". & \\
 & 2. Hiển thị thông tin ví hiện tại (nếu có). \\
3. Chọn "Kết nối ví". & \\
4. Đăng nhập vào ví Metamask và xác nhận kết nối. & \\
 & 5. Kiểm tra tính hợp lệ của địa chỉ ví (định dạng Ethereum). \\
 & 6. Cập nhật địa chỉ ví của học viên vào cơ sở dữ liệu và thông báo thành công. \\
\end{tabular}
\vspace{2pt}
\end{minipage} \\
\hline
\textbf{Luồng ngoại lệ} & 
\begin{minipage}[t]{\linewidth}
\begin{itemize}[nosep,after=\strut]
    \item \textbf{A1 -- Lỗi đăng nhập ví Metamask:} Thao tác bị hủy, hệ thống yêu cầu thử lại.
\end{itemize}
\end{minipage} \\
\hline
\textbf{Hậu điều kiện} & 
\begin{minipage}[t]{\linewidth}
\begin{itemize}[nosep,after=\strut]
    \item Địa chỉ ví của học viên được cập nhật thành công trong hệ thống.
    \item Các chứng chỉ NFT sau này sẽ được gửi đến địa chỉ ví mới.
\end{itemize}
\end{minipage} \\
\hline
\textbf{Quy tắc nghiệp vụ} & 
\begin{minipage}[t]{\linewidth}
\begin{itemize}[nosep,after=\strut]
    \item Địa chỉ ví phải tuân theo định dạng chuẩn Ethereum (bắt đầu bằng "0x" và gồm 42 ký tự).
\end{itemize}
\end{minipage} \\
\hline
\end{longtable}
Bảng \ref{tab:uc11} đặc tả ca sử dụng Cập nhật địa chỉ ví.

\subsubsection{Biểu đồ tuần tự}
\begin{figure}[H]
    \centering
    \includegraphics[width=0.8\textwidth]{img/wallet-se.png}
    \caption{Biểu đồ tuần tự ca sử dụng Cập nhật địa chỉ ví}
    \label{fig:uc11-se}
\end{figure}
Hình \ref{fig:uc11-se} trình bày về biểu đồ tuần tự ca sử dụng Cập nhật địa chỉ ví.

\subsubsection{Biểu đồ hoạt động}
\begin{figure}[H]
    \centering
    \includegraphics[width=0.8\textwidth]{img/wallet-ac.png}
    \caption{Biểu đồ hoạt động ca sử dụng Cập nhật địa chỉ ví}
    \label{fig:uc11-ac}
\end{figure}
Hình \ref{fig:uc11-ac} trình bày về biểu đồ hoạt động ca sử dụng Cập nhật địa chỉ ví.


%%%% UC-12: Nhận chứng chỉ %%%%
\subsection{Nhận chứng chỉ}
\subsubsection{Mô tả}
\begin{longtable}{|p{2.2cm}|p{13cm}|}
\caption{Đặc tả ca sử dụng Nhận chứng chỉ}
\label{tab:uc12} \\
\hline
\textbf{ID} & UC-12 \\
\hline
\textbf{Tên Use Case} & Nhận chứng chỉ \\
\hline
\textbf{Mô tả} & Học viên ký nhận chứng chỉ dưới dạng NFT sau khi được Đơn vị đào tạo cấp trên hệ thống. \\
\hline
\textbf{Actor} & Học viên \\
\hline
\textbf{Tiền điều kiện} & 
\begin{minipage}[t]{\linewidth}
\begin{itemize}[nosep,after=\strut]
    \item Học viên đã đăng nhập vào hệ thống.
    \item Học viên đã hoàn thành khóa học và đủ điều kiện nhận chứng chỉ.
    \item Đơn vị đào tạo đã cấp chứng chỉ NFT cho học viên.
    \item Học viên đã cập nhật địa chỉ ví hợp lệ vào hệ thống.
    \item Ví Metamask của học viên có ít nhất 0.01 SepoliaETH.
\end{itemize}
\end{minipage} \\
\hline
\textbf{Luồng chính} & 
\begin{minipage}[t]{\linewidth}
\vspace{2pt}
\begin{tabular}{p{6cm}|p{6cm}}
\textbf{Actor} & \textbf{System} \\
\hline
1. Truy cập trang "Chứng chỉ của tôi". & \\
 & 2. Hiển thị danh sách chứng chỉ NFT mà học viên đã được cấp. \\
3. Chọn chứng chỉ cần nhận. & \\
 & 4. Hiển thị thông tin chi tiết chứng chỉ. \\
5. Chọn "Nhận chứng chỉ" và xác nhận (Confirm) giao dịch trên Metamask. & \\
 & 6. Gửi giao dịch lên blockchain để mint NFT cho học viên. \\
 & 7. Cập nhật trạng thái chứng chỉ thành "Active" và thông báo thành công. \\
\end{tabular}
\vspace{2pt}
\end{minipage} \\
\hline
\textbf{Luồng ngoại lệ} & 
\begin{minipage}[t]{\linewidth}
\begin{itemize}[nosep,after=\strut]
    \item \textbf{A1 -- Lỗi giao dịch trên blockchain (không đủ số dư):} Hệ thống hiển thị trạng thái thất bại và yêu cầu nạp thêm ETH.
\end{itemize}
\end{minipage} \\
\hline
\textbf{Hậu điều kiện} & 
\begin{minipage}[t]{\linewidth}
\begin{itemize}[nosep,after=\strut]
    \item Chứng chỉ NFT được ghi nhận và liên kết chính thức với địa chỉ ví của học viên.
    \item Học viên có thể kiểm tra chứng chỉ trực tiếp trên blockchain.
\end{itemize}
\end{minipage} \\
\hline
\textbf{Quy tắc nghiệp vụ} & 
\begin{minipage}[t]{\linewidth}
\begin{itemize}[nosep,after=\strut]
    \item Chứng chỉ NFT là duy nhất (non-fungible) và không thể chuyển nhượng (Soulbound Token).
\end{itemize}
\end{minipage} \\
\hline
\end{longtable}
Bảng \ref{tab:uc12} đặc tả ca sử dụng Nhận chứng chỉ.

\subsubsection{Biểu đồ tuần tự}
\begin{figure}[H]
    \centering
    \includegraphics[width=0.8\textwidth]{img/claim-se.png}
    \caption{Biểu đồ tuần tự ca sử dụng Nhận chứng chỉ}
    \label{fig:uc12-se}
\end{figure}
Hình \ref{fig:uc12-se} trình bày về biểu đồ tuần tự ca sử dụng Nhận chứng chỉ.

\subsubsection{Biểu đồ hoạt động}
\begin{figure}[H]
    \centering
    \includegraphics[width=0.8\textwidth]{img/claim-ac.png}
    \caption{Biểu đồ hoạt động ca sử dụng Nhận chứng chỉ}
    \label{fig:uc12-ac}
\end{figure}
Hình \ref{fig:uc12-ac} trình bày về biểu đồ hoạt động ca sử dụng Nhận chứng chỉ.


%%%% UC-13: Xác minh chứng chỉ %%%%
\subsection{Xác minh chứng chỉ}
\subsubsection{Mô tả}
\begin{longtable}{|p{2.2cm}|p{13cm}|}
\caption{Đặc tả ca sử dụng Xác minh chứng chỉ}
\label{tab:uc13} \\
\hline
\textbf{ID} & UC-13 \\
\hline
\textbf{Tên Use Case} & Xác minh chứng chỉ \\
\hline
\textbf{Mô tả} & Nhà tuyển dụng tra cứu chứng chỉ của học viên và xác minh tính hợp lệ của chứng chỉ thông qua hệ thống và dữ liệu blockchain. \\
\hline
\textbf{Actor} & Nhà tuyển dụng \\
\hline
\textbf{Tiền điều kiện} & 
\begin{minipage}[t]{\linewidth}
\begin{itemize}[nosep,after=\strut]
    \item Nhà tuyển dụng có mã xác minh hoặc Token ID của chứng chỉ.
\end{itemize}
\end{minipage} \\
\hline
\textbf{Luồng chính} & 
\begin{minipage}[t]{\linewidth}
\vspace{2pt}
\begin{tabular}{p{6cm}|p{6cm}}
\textbf{Actor} & \textbf{System} \\
\hline
1. Truy cập trang xác minh chứng chỉ. & \\
 & 2. Hiển thị giao diện xác minh. \\
3. Nhập Token ID hoặc mã xác minh của chứng chỉ NFT. & \\
 & 4. Truy vấn thông tin chứng chỉ trên blockchain. \\
 & 5. Hiển thị thông tin chứng chỉ: tên học viên, tên khóa học, ngày cấp, Token ID, trạng thái (Active / Revoked / Replaced). \\
6. Xác nhận tính chính xác và hợp lệ của chứng chỉ. & \\
\end{tabular}
\vspace{2pt}
\end{minipage} \\
\hline
\textbf{Luồng ngoại lệ} & 
\begin{minipage}[t]{\linewidth}
\begin{itemize}[nosep,after=\strut]
    \item \textbf{A1 -- Không tìm thấy chứng chỉ:} Hệ thống hiển thị thông báo "Chứng chỉ không tồn tại".
\end{itemize}
\end{minipage} \\
\hline
\textbf{Hậu điều kiện} & 
\begin{minipage}[t]{\linewidth}
\begin{itemize}[nosep,after=\strut]
    \item Nhà tuyển dụng xác minh được thông tin chứng chỉ của học viên.
\end{itemize}
\end{minipage} \\
\hline
\textbf{Quy tắc nghiệp vụ} & 
\begin{minipage}[t]{\linewidth}
\begin{itemize}[nosep,after=\strut]
    \item Tất cả thông tin chứng chỉ phải được xác minh trực tiếp từ blockchain.
    \item Nếu chứng chỉ bị thu hồi hoặc thay thế, hệ thống phải hiển thị rõ ràng.
    \item Nhà tuyển dụng không cần tài khoản vẫn có quyền xác minh chứng chỉ.
    \item Dữ liệu chứng chỉ phải được lấy từ smart contract chính thức của hệ thống.
\end{itemize}
\end{minipage} \\
\hline
\end{longtable}
Bảng \ref{tab:uc13} đặc tả ca sử dụng Xác minh chứng chỉ.

\subsubsection{Biểu đồ tuần tự}
\begin{figure}[H]
    \centering
    \includegraphics[width=0.8\textwidth]{img/verify-se.png}
    \caption{Biểu đồ tuần tự ca sử dụng Xác minh chứng chỉ}
    \label{fig:uc13-se}
\end{figure}
Hình \ref{fig:uc13-se} trình bày về biểu đồ tuần tự ca sử dụng Xác minh chứng chỉ.

\subsubsection{Biểu đồ hoạt động}
\begin{figure}[H]
    \centering
    \includegraphics[width=0.8\textwidth]{img/verify-ac.png}
    \caption{Biểu đồ hoạt động ca sử dụng Xác minh chứng chỉ}
    \label{fig:uc13-ac}
\end{figure}
Hình \ref{fig:uc13-ac} trình bày về biểu đồ hoạt động ca sử dụng Xác minh chứng chỉ.


%%%% UC-14: Tóm tắt chứng chỉ bằng AI %%%%
\subsection{Tóm tắt chứng chỉ bằng AI}
\subsubsection{Mô tả}
\begin{longtable}{|p{2.2cm}|p{13cm}|}
\caption{Đặc tả ca sử dụng Tóm tắt chứng chỉ bằng AI}
\label{tab:uc14} \\
\hline
\textbf{ID} & UC-14 \\
\hline
\textbf{Tên Use Case} & Tóm tắt chứng chỉ bằng AI \\
\hline
\textbf{Mô tả} & Hệ thống AI phân tích và tóm tắt chứng chỉ để nhà tuyển dụng hiểu nhanh về nội dung và kỹ năng của ứng viên. \\
\hline
\textbf{Actor} & Nhà tuyển dụng \\
\hline
\textbf{Tiền điều kiện} & 
\begin{minipage}[t]{\linewidth}
\begin{itemize}[nosep,after=\strut]
    \item Nhà tuyển dụng đã truy cập trang xác minh chứng chỉ.
    \item Chứng chỉ đã được xác minh thành công và thông tin chi tiết đang được hiển thị.
\end{itemize}
\end{minipage} \\
\hline
\textbf{Luồng chính} & 
\begin{minipage}[t]{\linewidth}
\vspace{2pt}
\begin{tabular}{p{6cm}|p{6cm}}
\textbf{Actor} & \textbf{System} \\
\hline
1. Tại trang chi tiết chứng chỉ, chọn nút ``Tóm tắt chứng chỉ bằng AI''. & \\
 & 2. Hiển thị hộp thoại tóm tắt AI với trạng thái đang phân tích. \\
 & 3. Gửi thông tin khóa học (tên, thời lượng, mô tả, nội dung đào tạo) đến LLM. \\
 & 4. Nhận kết quả từ LLM và hiển thị bản tóm tắt AI cho nhà tuyển dụng. \\
5. Đọc bản tóm tắt để hiểu nhanh nội dung và kỹ năng của ứng viên. & \\
\end{tabular}
\vspace{2pt}
\end{minipage} \\
\hline
\textbf{Luồng ngoại lệ} & 
\begin{minipage}[t]{\linewidth}
\begin{itemize}[nosep,after=\strut]
    \item \textbf{A1 -- Lỗi kết nối đến dịch vụ AI:} Hệ thống hiển thị thông báo lỗi và yêu cầu thử lại.
\end{itemize}
\end{minipage} \\
\hline
\textbf{Hậu điều kiện} & 
\begin{minipage}[t]{\linewidth}
\begin{itemize}[nosep,after=\strut]
    \item Nhà tuyển dụng nhận được bản tóm tắt ngắn gọn về chứng chỉ và kỹ năng của ứng viên.
\end{itemize}
\end{minipage} \\
\hline
\textbf{Quy tắc nghiệp vụ} & 
\begin{minipage}[t]{\linewidth}
\begin{itemize}[nosep,after=\strut]
    \item Bản tóm tắt được tạo bởi mô hình AI (GPT) dựa trên dữ liệu khóa học thực tế.
    \item Nhà tuyển dụng có thể yêu cầu tạo lại tóm tắt nếu chưa hài lòng.
\end{itemize}
\end{minipage} \\
\hline
\end{longtable}
Bảng \ref{tab:uc14} đặc tả ca sử dụng Tóm tắt chứng chỉ bằng AI.

\subsubsection{Biểu đồ tuần tự}
\begin{figure}[H]
    \centering
    \includegraphics[width=0.8\textwidth]{img/AI-summary-se.png}
    \caption{Biểu đồ tuần tự ca sử dụng Tóm tắt chứng chỉ bằng AI}
    \label{fig:uc14-se}
\end{figure}
Hình \ref{fig:uc14-se} trình bày về biểu đồ tuần tự ca sử dụng Tóm tắt chứng chỉ bằng AI.

\subsubsection{Biểu đồ hoạt động}
\begin{figure}[H]
    \centering
    \includegraphics[width=0.6\textwidth]{img/AI-summary-ac.png}
    \caption{Biểu đồ hoạt động ca sử dụng Tóm tắt chứng chỉ bằng AI}
    \label{fig:uc14-ac}
\end{figure}
Hình \ref{fig:uc14-ac} trình bày về biểu đồ hoạt động ca sử dụng Tóm tắt chứng chỉ bằng AI.


%%%% UC-15: Tương tác chatbot %%%%
\subsection{Tương tác chatbot}
\subsubsection{Mô tả}
\begin{longtable}{|p{2.2cm}|p{13cm}|}
\caption{Đặc tả ca sử dụng Tương tác chatbot}
\label{tab:uc15} \\
\hline
\textbf{ID} & UC-15 \\
\hline
\textbf{Tên Use Case} & Tương tác chatbot \\
\hline
\textbf{Mô tả} & Người dùng tương tác với chatbot AI để hỏi về cách sử dụng hệ thống, nội dung chương trình đào tạo và các vấn đề liên quan đến chứng chỉ blockchain. \\
\hline
\textbf{Actor} & Đơn vị đào tạo, Học viên, Nhà tuyển dụng \\
\hline
\textbf{Tiền điều kiện} & 
\begin{minipage}[t]{\linewidth}
\begin{itemize}[nosep,after=\strut]
    \item Người dùng đã đăng nhập vào hệ thống.
\end{itemize}
\end{minipage} \\
\hline
\textbf{Luồng chính} & 
\begin{minipage}[t]{\linewidth}
\vspace{2pt}
\begin{tabular}{p{6cm}|p{6cm}}
\textbf{Actor} & \textbf{System} \\
\hline
1. Nhấn vào biểu tượng chatbot ở góc màn hình. & \\
 & 2. Mở cửa sổ chatbot và hiển thị lời chào cùng danh sách gợi ý câu hỏi. \\
3. Nhập câu hỏi hoặc chọn một câu hỏi gợi ý. & \\
 & 4. Phân tích nội dung câu hỏi và truy xuất thông tin liên quan từ cơ sở tri thức. \\
 & 5. Hiển thị câu trả lời cho người dùng. \\
6. Tiếp tục đặt câu hỏi hoặc đóng cửa sổ chatbot. & \\
\end{tabular}
\vspace{2pt}
\end{minipage} \\
\hline
\textbf{Luồng ngoại lệ} & 
\begin{minipage}[t]{\linewidth}
\begin{itemize}[nosep,after=\strut]
    \item \textbf{A1 -- Câu hỏi nằm ngoài phạm vi hỗ trợ:} Chatbot phản hồi gợi ý câu hỏi phù hợp hoặc hướng dẫn liên hệ hỗ trợ kỹ thuật.
    \item \textbf{A2 -- Lỗi kết nối đến dịch vụ AI:} Hệ thống thông báo không thể xử lý yêu cầu và yêu cầu thử lại.
\end{itemize}
\end{minipage} \\
\hline
\textbf{Hậu điều kiện} & 
\begin{minipage}[t]{\linewidth}
\begin{itemize}[nosep,after=\strut]
    \item Người dùng nhận được câu trả lời cho câu hỏi của mình.
    \item Lịch sử hội thoại được duy trì trong phiên làm việc.
\end{itemize}
\end{minipage} \\
\hline
\textbf{Quy tắc nghiệp vụ} & 
\begin{minipage}[t]{\linewidth}
\begin{itemize}[nosep,after=\strut]
    \item Chatbot hỗ trợ về cách sử dụng hệ thống hoặc nội dung chương trình đào tạo.
    \item Gợi ý câu hỏi nhanh được hiển thị sẵn để hỗ trợ người dùng.
\end{itemize}
\end{minipage} \\
\hline
\end{longtable}
Bảng \ref{tab:uc15} đặc tả ca sử dụng Tương tác chatbot.

\subsubsection{Biểu đồ tuần tự}
\begin{figure}[H]
    \centering
    \includegraphics[width=0.8\textwidth]{img/chatbot-se.png}
    \caption{Biểu đồ tuần tự ca sử dụng Tương tác chatbot}
    \label{fig:uc15-se}
\end{figure}
Hình \ref{fig:uc15-se} trình bày về biểu đồ tuần tự ca sử dụng Tương tác chatbot.

\subsubsection{Biểu đồ hoạt động}
\begin{figure}[H]
    \centering
    \includegraphics[width=0.8\textwidth]{img/chatbot-ac.png}
    \caption{Biểu đồ hoạt động ca sử dụng Tương tác chatbot}
    \label{fig:uc15-ac}
\end{figure}
Hình \ref{fig:uc15-ac} trình bày về biểu đồ hoạt động ca sử dụng Tương tác chatbot.
